% ---------------------------------------------------------------------
%  Chapter 14 : The Spatial Integrator --- Symanzik Polynomials and
%               Causal Chambers
% ---------------------------------------------------------------------
\chapter{The Spatial Integrator: Symanzik Polynomials and Causal Chambers}\label{ch:spatial-core}

The previous part of this manual evaluated diagrams \emph{in the time domain}:
each diagram became a product of exponential propagators, and the loop integrals
were done over the internal vertex times. That route works beautifully for
temporal (time-only) models, but it carries a structural wound --- the
\term{close-pair pathology}, a $1/\beta$ factor that blows up whenever two poles
of the time integrand approach each other. For a \emph{spatial} model --- a
stochastic \emph{partial} differential equation, where the field lives on a
$d$-dimensional space as well as in time --- there is a much cleaner road, and
this chapter is about the subsystem that paves it.

The idea, in one sentence, is: \emph{do the momentum integral first, and do it
analytically}. Every enumerated diagram --- a tree, a one-loop bubble, a one-loop
tadpole, a two-loop sunset, for any number of external legs $k$, any loop order
$\ell$, and any spatial dimension $d$ --- is reduced to the \emph{same} integral
and evaluated by the \emph{same} code. There is no per-topology branch. The
phrase the codebase repeats, in the very first lines of the production module, is
that ``ONE genuine integral evaluates \emph{every} enumerated diagram''
(\file{msrjd/integration/spatial/full\_integrator.py:4}). This chapter explains
why that is possible and walks the machinery that does it.

\begin{note}
This is the first of three spatial chapters. This one builds the core ---
momentum routing, the Symanzik polynomials, the causal chambers, and the kinematic
integral. Chapter~\ref{ch:spatial-heatkernel} covers the analytic
inverse Fourier transform and the derivative (form-factor) vertices in detail.
Chapter~\ref{ch:spatial-coupled} covers the coupled multi-field (unequal-$D$)
extension. Everything here lives under
\file{msrjd/integration/spatial/}, and none of it imports SageMath or
\code{nauty} --- those are upstream, in the enumeration. The only heavy
dependencies here are NumPy, SciPy, and (in one small place) SymPy.
\end{note}

\section{Why momentum-first}

\subsection{The wound in the temporal route}

Recall how the temporal pipeline evaluates a loop. After enumerating and typing a
diagram it has a product of propagators in the $(\,$frequency, or time$\,)$
representation, and it integrates the internal vertex times by carving the time
domain into ordering regions and integrating a product of exponentials over each.
The integrand of one such region looks schematically like a sum of terms
$\ee^{-\beta_1 t_1}\ee^{-\beta_2 t_2}\cdots$, and doing the nested integrals
produces denominators built from \emph{differences} of the decay rates $\beta_i$.
When two of those rates coincide --- a degenerate or near-degenerate pole pair ---
the denominator $1/(\beta_i-\beta_j)$ diverges. The true integral is finite (the
apparent pole is removable, a confluent limit), but the naive expression is not,
and resolving it costs precision and special-case code. That is the close-pair
pathology, and Part~IV spends real effort taming it.

\subsection{The cure}

For a spatial model the field also carries a momentum. The crucial observation is
that in the heat-kernel representation the edge \emph{times are already Schwinger
parameters} (we make this precise in \S\ref{sec:heatkernel} below), so the
loop-\emph{momentum} integral $\int \dd^d\ell$ is a \emph{pure Gaussian}. A
Gaussian integral can be done in closed form, exactly, at any loop order and in
any dimension --- it collapses to the two \term{Symanzik polynomials}. And here is
the payoff: the close-pair pathology lived entirely in the \emph{momentum}
integral. Once that integral is done analytically, what remains is a smooth
integral over the internal vertex times and a few auxiliary parameters, with no
poles to collide. Ordinary quadrature converges fast, and close-pair
\emph{cannot arise}. The comment block in \file{causal\_chambers.py:17-24} states
this explicitly: the temporal pipeline ``integrates products of exponentials per
chamber $\dots$ whose $1/\beta$ factors are the close-pair pathology. Here the
loop momenta are already integrated $\dots$ so the per-chamber integrand is
\emph{smooth} $\dots$ and close-pair cannot arise.''

\begin{defn}
The \term{one mental model} for this whole subsystem: a heat-kernel propagator in
the momentum--time representation is $G(k,t)=\theta(t)\,\ee^{-(\mu+Dk^2)t}$. The
edge times act as Schwinger parameters, so $\int\dd^d\ell$ is a pure Gaussian and
collapses \emph{exactly} to the Symanzik polynomials $\Usym$ and $\Fsym$ at any
loop order $L$ and any dimension $d$. After that collapse the only thing left is a
smooth integral over the internal vertex times and the noise-source Schwinger
parameters --- no pole, no close pair. That is the entire design.
\end{defn}

\subsection{Where this sits in the pipeline}

The subsystem consumes a fully enumerated and typed diagram and produces that
diagram's contribution to the connected $k$-point cumulant. Upstream, the
enumeration has (a)~expanded the action into vertices, (b)~enumerated all
topologically distinct Feynman diagrams up to the chosen loop order, and
(c)~``typed'' each diagram --- decided which lines are response/retarded ($R$)
lines and which are noise/correlation ($C$) lines, and attached a symbolic
prefactor carrying the couplings, noise amplitudes, and symmetry factor
$\Sym(\Gamma)$. What is left is a hard multidimensional integral \emph{per
diagram}, and this subsystem performs it. The flow is:

\begin{lstlisting}[style=console]
   theory (.theory.py)  ->  action -> vertices
        |
   diagram enumeration  (msrjd/diagrams/...)   -> prediagrams
        |  field-type assignment
        v
   TypedDiagram          (msrjd/diagrams/type_assignment.py)
        |  .prediagram, .vertex_assignments, .propagator_indices, prefactor
        v
 +-----------------------  THIS SUBSYSTEM  --------------------------+
 |  diagram_to_cstack(td)              -> CStackDiagram (descriptor) |
 |  route_momenta(td) -> edge_coeffs() -> per-edge (a_e, b_e)        |
 |  _symanzik_kernel_batch / _momentum_factor_batch -> U, F          |
 |  causal_chambers(...)               -> orderings of {t_v}         |
 |  diagram_kinematic(...)             -> int dt int dsigma          |
 |  _heat_kernel_x_general / IFT       -> dC(x,tau)  directly        |
 +------------------------------------------------------------------+
        |  sum over diagrams, x 2^{-n_C} . S(Gamma) . prefactor
        v
   C(q,tau) or dC(x,tau)  ->  pipeline_bridge.py  ->  compute_cumulants(...)
\end{lstlisting}

The six files of the subsystem, and their roles:

\begin{center}
\begin{tabular}{@{}>{\ttfamily\footnotesize}l l@{}}
\toprule
\normalfont\textbf{file (under \file{msrjd/integration/spatial/})} & \textbf{role} \\
\midrule
diagram\_descriptor.py & typed diagram $\to$ flat ``C-stack'' edge list \\
momentum\_routing.py   & momentum conservation $\to$ routing coefficients $(a_e,b_e)$ \\
spatial\_reduce.py     & the textbook Symanzik reduction (single-sample reference) \\
causal\_chambers.py    & the causal time-ordering chambers \\
full\_integrator.py    & the production integrator (batched + analytic IFT) \\
generic\_evaluator.py  & \normalfont\emph{oracle only} --- an older cross-check \\
\bottomrule
\end{tabular}
\end{center}

\section{The heat-kernel representation}\label{sec:heatkernel}

This section builds the relevant theory from the ground up. A reader who knows
\msrjd{} field theory but not this particular representation should be able to
follow every step.

\subsection{The two propagators}

Start from a linearized (Gaussian, ``tree'') theory with one field. In the
momentum--time representation the \term{retarded propagator} --- the response
function, the $R$ line --- is
\begin{equation}
  \Gret(k,t) \;=\; \theta(t)\,\ee^{-m(k)\,t},
  \qquad
  m(k) \;=\; \mu + D\,k^2,
  \label{eq:GR}
\end{equation}
where $\mu$ is a mass / relaxation rate, $D$ a diffusion constant,
$k^2=|k|^2$ the squared momentum, and $\theta$ the Heaviside step. The step
function is causality made manifest: a response only to the past. The
\term{correlation function} --- the $C$ line, $\avg{\phi\phi}$ --- of the
\emph{tree} theory is
\begin{equation}
  C_0(k,t) \;=\; \frac{T}{m(k)}\;\ee^{-m(k)\,|t|},
  \label{eq:C0}
\end{equation}
with $T$ the noise strength (the action carries a noise vertex $-T\,\tilde\phi^2$,
or $2T\,\tilde\phi\,\tilde\phi$ in the symmetric normalization). Real-space
correlators are inverse Fourier transforms of these over $k$.

\subsection{A correlation line is two retarded segments glued at a noise source}

The identity the whole subsystem rests on is that a correlation line can be
written as two retarded segments joined at a noise source, with the source time
integrated out. Schematically,
\begin{equation}
  C_0(\Delta t) \;=\; \frac{T}{m}\,\ee^{-m|\Delta t|}
  \;=\; T\!\int_0^{\infty}\!\dd\sigma\;\ee^{-m(|\Delta t|+\sigma)}\cdot(\text{const}),
  \label{eq:Cglue}
\end{equation}
i.e.\ the $C$ line is an $R$--noise--$R$ object. Integrating the noise-source
time produces a single \term{Schwinger parameter} $\sigma\in[0,\infty)$ and the
weight $\ee^{-m(|\Delta t|+\sigma)}$; the convention factor of $2$ is bookkept by
$2^{-n_C}$ (\S\ref{sec:norm}). This is exactly why the enumerator can work
entirely with all-$\Gret$ propagators plus explicit noise sources, and why the
descriptor re-contracts a 2-point noise source back into one $C$ edge.

\begin{defn}
A \term{Schwinger parameter} is an auxiliary integration variable that linearizes
a propagator denominator into an exponential. Here there are two kinds: an
$R$-edge time $w$, and a noise-source parameter $\sigma\in[0,\infty)$. Their job
is to turn the momentum dependence of every propagator into a single Gaussian
exponent $\exp(-D\sum_e w_e k_e^2)$, so that $\int\dd^d\ell$ is exactly doable.
\end{defn}

\subsection{Edge weights and the loop-momentum integral}

For a diagram with internal vertices at times $\{t_v\}$ (to be integrated),
external leaves at fixed times $\{\tau_j\}$, and edges $e$, each edge carries a
\term{Schwinger weight} $w_e$ (\file{full\_integrator.py:16-19}):
\begin{align}
  R \text{ edge } (t_{\text{head}},t_{\text{tail}}):&\quad
     w_e = t_{\text{head}} - t_{\text{tail}} \;\ge\; 0 \ \text{(by causality)},\\
  C \text{ edge between } t_a,t_b:&\quad
     w_e = |t_a - t_b| + \sigma_e,\quad \sigma_e\in[0,\infty).
\end{align}
Each edge also carries a routed momentum $k_e$ (next section). The diagram's
loop-momentum integral, in the heat-kernel representation, is the pure Gaussian
\begin{equation}
  I_{\text{mom}} \;=\; \int \prod_{i=1}^{L}\frac{\dd^d\ell_i}{(2\pi)^d}\;
     \exp\!\Big[-D\sum_e w_e\,k_e^2\Big],
  \label{eq:Imom}
\end{equation}
with $L$ the number of loops. This is the central object
(\file{spatial\_reduce.py:12}). The next two sections compute it.

\section{Momentum routing}\label{sec:routing}

\subsection{Conservation at a spatial vertex}

A spatial interaction vertex is \emph{local in space}: the action contains
$\int\dd x_v$ of the product of the fields meeting at the vertex. Fourier
transforming, that integral enforces \term{momentum conservation}
$\sum_e k_e = 0$ at every internal vertex --- the spatial analogue of the
time-domain vertex, which carries one integrated time $t_v$. Solving this linear
conservation system writes each edge momentum as a linear form
\begin{equation}
  k_e \;=\; \sum_{i=1}^{L} a_{ei}\,\ell_i \;+\; \sum_{j=1}^{n_{\text{ext}}} b_{ej}\,q_j,
  \label{eq:route}
\end{equation}
in the \term{loop momenta} $\ell_i$ (the $L$ free parameters of the conservation
system) and the \term{external momenta} $q_j$ (one per external leg). Overall
conservation $\sum_j q_j = 0$ is imposed by eliminating the last leaf, so
$n_{\text{ext}} = k-1$. The coefficients $(a_e,b_e)$ are the per-edge
\term{routing coefficients}; for a translation-invariant routing they are
integers, $\pm1$ or $0$.

Two cases make this concrete (\file{momentum\_routing.py:29-34}):

\begin{itemize}
  \item \textbf{Tree ($L=0$).} Every edge carries $\pm q$ only, so $k_e^2 = q^2$
        uniformly. A single-mode tree therefore needs only a
        $\Lap \to -q^2$ substitution.
  \item \textbf{One-loop bubble.} The two internal lines carry $\ell$ and
        $q-\ell$ respectively --- genuinely different momenta --- so the
        substitution becomes per-edge.
\end{itemize}

\subsection{The symbolic solve}

The routing is done with SymPy, the pure-Python computer-algebra system, because
momentum conservation is a small \emph{exact} linear-algebra problem and a CAS
gives integer $\pm1/0$ coefficients with no floating-point ambiguity.

\begin{note}
\textbf{What SymPy is.} SymPy (\code{import sympy as sp}) represents mathematical
\emph{symbols} (\code{sp.Symbol}), does exact rational/integer arithmetic, and
--- crucially here --- solves linear systems symbolically with
\code{sp.linsolve}. It is the only place in this entire subsystem that SymPy
appears; the heavy numerics are all NumPy/SciPy.
\end{note}

The function is \code{route\_momenta} (\file{momentum\_routing.py:99}). Reading
its body, the steps are:

\begin{enumerate}
  \item Create the external-momentum symbols $q_0\dots q_{k-1}$ and one unknown
        momentum symbol per edge:
        \begin{lstlisting}
q = sp.symbols(f'q0:{k}') if k > 0 else ()
kk = {e: sp.Symbol(f'k{i}') for i, e in enumerate(edges)}   # one per edge
        \end{lstlisting}
  \item Build the conservation \term{residual} at each vertex --- the signed sum
        that must vanish. The convention (\file{momentum\_routing.py:103-107}) is
        that edge $(u,v,\text{lbl})$ is oriented $u\to v$, so it contributes
        $+k_e$ at the vertex it flows \emph{into} and $-k_e$ at the vertex it
        flows \emph{out of}; at a leaf the injected external momentum $q_j$ is
        subtracted:
        \begin{lstlisting}
for v in vertices:
    s = sp.Integer(0)
    for e in edges:
        u, w, _lbl = e
        if w == v: s += kk[e]          # edge flows INTO v
        if u == v: s -= kk[e]          # edge flows OUT of v
    if v in leaves:
        s -= q[leaves.index(v)]        # external momentum injected at leaf
    residuals.append(sp.expand(s))
        \end{lstlisting}
  \item Impose overall conservation by substituting
        $q_{k-1} = -(q_0+\dots+q_{k-2})$ into every residual.
  \item Solve the system for the edge momenta with SymPy's linear solver:
        \begin{lstlisting}
sol_set = sp.linsolve(residuals, list(kk.values()))
        \end{lstlisting}
        An empty solution set means the system is inconsistent --- which never
        happens for a valid \msrjd{} diagram, so the code raises a clear error if
        it does (\file{momentum\_routing.py:139-142}).
  \item The \emph{free symbols} SymPy leaves unpinned in the solution \emph{are}
        the loop momenta. The code collects them
        (\code{expr.free\_symbols}), subtracts the external symbols, sorts by
        name, and relabels them to $\ell_0\dots\ell_{L-1}$
        (\file{momentum\_routing.py:149-158}).
\end{enumerate}

\subsection{The routing result}

The result is a small dataclass (\file{momentum\_routing.py:50}):

\begin{lstlisting}
@dataclass
class RoutingResult:
    edge_momenta: dict   # (u,v,lbl) -> sympy expr in q_syms + loop_syms
    q_syms: tuple        # external-momentum symbols q_0 ... q_{k-1}
    loop_syms: tuple     # loop-momentum symbols l_0 ... l_{L-1}
    n_loops: int
\end{lstlisting}

It exposes two extractors. \code{edge\_coeffs()}
(\file{momentum\_routing.py:62}) returns \code{\{edge: (a, b)\}} --- the signed
linear coefficients of \eqref{eq:route}, each coerced to a Python \code{int} by a
small helper \code{\_num} (\file{momentum\_routing.py:79}). \textbf{This is what
the Symanzik step consumes.} The other extractor, \code{edge\_k2()}
(\file{momentum\_routing.py:57}), returns the symbolic $k_e^2$ --- but it must
\emph{not} be used for the reduction.

\begin{gotcha}
\code{edge\_k2()} squares the momentum and loses the \emph{signed} cross-term
coefficients (the $a_{ei}b_{ej}$ products that build the loop/external coupling
block). Only \code{edge\_coeffs()} carries the $(a,b)$ the Symanzik step needs.
The code warns about this in the docstring at
\file{momentum\_routing.py:72-74}. If you ever extend the integrator, consume
\code{edge\_coeffs()}, never \code{edge\_k2()}.
\end{gotcha}

\section{Symanzik polynomials from scratch}\label{sec:symanzik}

Now do the Gaussian integral \eqref{eq:Imom}. Substitute the routing
\eqref{eq:route} and collect the loop momenta into a quadratic form. Define three
matrices by summing the routing coefficients weighted by the Schwinger weights
(\file{spatial\_reduce.py:18-21}):
\begin{align}
  \Lambda_{ii'} &= \sum_e w_e\,a_{ei}\,a_{ei'} &&(L\times L) &&\text{first-Symanzik kernel,}\\
  N_{ij} &= \sum_e w_e\,a_{ei}\,b_{ej} &&(L\times n_{\text{ext}}) &&\text{loop/external cross block,}\\
  Q_{jj'} &= \sum_e w_e\,b_{ej}\,b_{ej'} &&(n_{\text{ext}}\times n_{\text{ext}}) &&\text{pure external block.}
\end{align}
Then $\sum_e w_e k_e^2 = \ell^{\mathsf T}\Lambda\ell + 2q^{\mathsf T}N^{\mathsf T}\ell
+ q^{\mathsf T}Q\,q$. Completing the square in $\ell$ and doing the Gaussian
integral gives the closed form (\file{spatial\_reduce.py:24-26}):
\begin{equation}
  I_{\text{mom}} \;=\; (4\pi D)^{-Ld/2}\;\Usym^{-d/2}\;
     \exp\!\big[-D\,q^{\mathsf T} Q_{\text{eff}}\, q\big],
  \label{eq:collapse}
\end{equation}
with
\begin{equation}
  \Usym = \det\Lambda,
  \qquad
  Q_{\text{eff}} = Q - N^{\mathsf T}\Lambda^{-1}N,
  \qquad
  \Fsym \equiv \Usym\cdot Q_{\text{eff}}\ \text{(scalar case)}.
  \label{eq:UF}
\end{equation}

\subsection{What these polynomials are}

It is worth saying what $\Usym$ and $Q_{\text{eff}}$ \emph{are} classically,
because a reader may know the graph-theoretic definitions.

\begin{description}
  \item[$\Usym$ --- the first Symanzik (first graph) polynomial.] A homogeneous
        degree-$L$ polynomial in the edge weights. Graph-theoretically,
        $\Usym = \sum_{\text{spanning trees }T}\ \prod_{e\notin T} w_e$ --- a sum
        over spanning trees of the product of the weights of the edges \emph{not}
        in the tree. In this code it is computed \emph{numerically} as
        $\det\Lambda$, which equals that combinatorial sum (the matrix-tree
        theorem). Validated examples (\file{spatial\_reduce.py:36-38}): the
        one-loop bubble has $\Usym = w_1+w_2$; the two-loop sunset has
        $\Usym = w_1w_2 + w_2w_3 + w_3w_1$.
  \item[$Q_{\text{eff}}$ --- the second Symanzik / $\Fsym$.] Carries the
        external-momentum dependence. $Q_{\text{eff}} = Q - N^{\mathsf T}\Lambda^{-1}N$
        is the \term{Schur complement} of $\Lambda$ in the full quadratic form ---
        exactly the residual external quadratic form after integrating out the
        loops. For the bubble, $Q_{\text{eff}} = w_1w_2/(w_1+w_2)$; for the
        sunset, $Q_{\text{eff}} = w_1w_2w_3/\Usym$. For $n_{\text{ext}}=1$ (the
        2-point function, $k=2$) $Q_{\text{eff}}$ is a scalar and
        $\Fsym = \Usym\cdot Q_{\text{eff}}$.
\end{description}

\begin{defn}
The \term{Schur complement} of a block $\Lambda$ in a partitioned quadratic form
$\left(\begin{smallmatrix}\Lambda & N\\ N^{\mathsf T} & Q\end{smallmatrix}\right)$ is
$Q - N^{\mathsf T}\Lambda^{-1}N$. It is precisely what is left of the $q$-quadratic
form after the $\ell$ block has been integrated out (Gaussian-eliminated). That is
why $Q_{\text{eff}}$ is the residual external form: integrating the loops
\emph{is} taking a Schur complement.
\end{defn}

The exponents $-Ld/2$ and $-d/2$ are the \emph{only} places the spatial dimension
$d$ enters the momentum reduction. The reduction is otherwise exact and
$d$-agnostic --- the same algebra works in $d=1$ as in $d=7$.

\subsection{The reference implementation}

\file{spatial\_reduce.py} is the textbook, single-sample reference. Its three
functions are deliberately small and quotable. \code{symanzik\_matrices}
(\file{spatial\_reduce.py:58}) builds $(\Lambda,N,Q)$ for one set of edge weights
by exactly the contractions above:
\begin{lstlisting}
aw  = a * w[:, None]          # (E, L)
Lam = aw.T @ a                # (L, L)
N   = aw.T @ b                # (L, n_ext)
Q   = (b * w[:, None]).T @ b  # (n_ext, n_ext)
\end{lstlisting}
\code{symanzik\_polynomials} (\file{spatial\_reduce.py:82}) returns
$(\Usym, Q_{\text{eff}})$ with $\Usym = \det\Lambda$ and $Q_{\text{eff}}$
computed via \code{np.linalg.solve} (not an explicit inverse --- solving a linear
system is both faster and more accurate than forming $\Lambda^{-1}$). And
\code{momentum\_integral} (\file{spatial\_reduce.py:101}) evaluates
\eqref{eq:collapse} for one sample. The production code re-implements all three,
\emph{batched}, in \file{full\_integrator.py}; the reference exists so the batched
version can be checked against it.

\begin{note}
\textbf{Why \code{np.linalg.solve} and not \code{np.linalg.inv}.} To compute
$N^{\mathsf T}\Lambda^{-1}N$ you never need $\Lambda^{-1}$ itself --- you need the
product $\Lambda^{-1}N$, which is the solution $X$ of the linear system
$\Lambda X = N$. \code{np.linalg.solve(Lam, N)} computes exactly that, using an
LU factorization rather than a full matrix inversion. It is the standard numerical
hygiene: never invert when you can solve.
\end{note}

\section{The full diagram integral}

Putting the pieces together, one diagram's contribution to the connected
$k$-point cumulant is (\file{full\_integrator.py:15-20}):
\begin{equation}
  \Gamma(q,\{\tau\}) \;=\; 2^{-n_C}\,\Sym(\Gamma)\!
     \int \prod_v \dd t_v \prod_{C\text{ edges}}\!\dd\sigma_e\;
     \mathbf{1}(\theta\text{'s})\;
     \ee^{-\mu\sum_e w_e}\;\text{MomFactor}(w,q),
  \label{eq:diagram}
\end{equation}
\begin{equation}
  \text{MomFactor} \;=\; (4\pi D)^{-Ld/2}\,\Usym(w)^{-d/2}\,
     \ee^{-D\,q^{\mathsf T}Q_{\text{eff}}(w)\,q}.
\end{equation}
The factor $\ee^{-\mu\sum_e w_e}$ is the \term{mass damping}: every segment of
length $w$ decays as $\ee^{-\mu w}$, since the $\mu$ part of the mass $m(k)=\mu+Dk^2$
multiplies $w_e$ directly. The $Dk^2$ part of the mass already lives \emph{inside}
$\text{MomFactor}$, through the Symanzik exponent. The $\mathbf{1}(\theta\text{'s})$
are the causal step functions of the retarded edges. The residual integral is over
the internal vertex times $t_v$ and the correlation Schwinger parameters
$\sigma_e$ only --- and it is \emph{smooth}, because the loop integral, the source
of the close-pair pathology, was already done analytically.

\section{Normalization --- derived, not fitted}\label{sec:norm}

The integrator's normalization is \emph{derived}, not tuned to a fixture
(\file{full\_integrator.py:27-31}). The enumeration prefactor uses the $2T$
noise-vertex convention. A kinematic $C$ edge here, by contrast, is the
unit-amplitude Schwinger factor $\int_0^\infty\dd\sigma\,\ee^{-m\sigma}=1/m$. The
factor $2^{-n_C}$ ($n_C$ = number of $C$ edges) converts between the two. The
tree check makes this airtight: a tree 2-point function has no loops and one $C$
edge between the two leaves ($n_C=1$), prefactor $2T$, and the integral gives
\begin{equation}
  2^{-1}\cdot 2T\cdot\frac{1}{m}\,\ee^{-m\tau}
  \;=\; \frac{T}{m}\,\ee^{-m\tau}
  \;=\; C_0,
\end{equation}
exactly the tree correlator \eqref{eq:C0}. This $2^{-n_C}$ is applied in
\code{diagram\_value} / \code{diagram\_value\_x} (\file{full\_integrator.py:1153},
\file{full\_integrator.py:1210}).

\section{Causal chambers}\label{sec:chambers}

\subsection{Posets and linear extensions}

What remains after the momentum collapse is
$\int\prod_v\dd t_v\,\mathbf{1}(\theta\text{'s})\,(\text{smooth integrand})$. The
retarded edges impose orderings $t_v > t_u$ on the internal vertices --- a
\term{partial order} (a poset). The domain of a poset decomposes into
\term{ordering chambers}, one per \term{linear extension} (a total order
consistent with the partial order). The key fact: within each chamber every
$|\Delta t|$ has a \emph{fixed sign}, so the integrand is smooth (no cusps) and
ordinary quadrature converges fast (\file{causal\_chambers.py:18-24}).

\begin{defn}
A \term{linear extension} of a poset is a total order of its elements that
respects all the partial-order relations. Geometrically, each linear extension is
one simplicial chamber of the cube $[\,\text{lo},\text{hi}\,]^{n}$, and the
chambers \emph{partition} the constrained domain. Integrating over the whole
constrained domain therefore equals summing the integral over the chambers --- no
over- or under-counting.
\end{defn}

\subsection{The code, reused verbatim from the temporal pipeline}

The poset machinery is not reimplemented for the spatial path. The single public
function (\file{causal\_chambers.py:41}) wraps the temporal pipeline's
battle-tested \code{\_CausalPoset} and \code{\_enumerate\_linear\_extensions}:

\begin{lstlisting}
def causal_chambers(n_vertices, retarded_edges):
    poset = _CausalPoset(
        m=int(n_vertices),
        edges=tuple(sorted({(int(u), int(v)) for (u, v) in retarded_edges})),
        scalar_lowers=(), scalar_uppers=())
    return list(_enumerate_linear_extensions(poset))
\end{lstlisting}

Here \code{retarded\_edges} is an iterable of pairs $(u,v)$ meaning $t_v > t_u$.
The function returns the list of linear extensions, each a length-$n$ tuple giving
the vertex order (earliest first).

\begin{gotcha}
\textbf{The empty-edge case tiles the cube.} With \emph{no} retarded ordering
constraints, \code{causal\_chambers} returns all $n!$ orderings, which
\emph{partition} the time cube --- the integral over the cube equals the sum over
chambers. Forgetting this gives an $n!$ over- or under-count. This is why the
production integrator always loops over the returned chambers even when there are
no ordering constraints.
\end{gotcha}

The module also ships a reference integrator, \code{integrate\_over\_chambers}
(\file{causal\_chambers.py:84}), which does the chamber simplex integral by nested
\code{scipy.integrate.quad} (adaptive Gauss--Kronrod 1-D quadrature). That path is
an oracle; the \emph{production} integrator uses fixed Gauss--Legendre product
grids instead, for vectorization, which we turn to next.

\section{The descriptor: \texorpdfstring{\texttt{diagram\_to\_cstack}}{diagram\_to\_cstack}}\label{sec:descriptor}

Before the integrator can run, the typed diagram must be flattened into the
canonical \term{C-stack} form that every diagram shares. That is the job of
\code{diagram\_to\_cstack} (\file{diagram\_descriptor.py:105}). The two record
types it produces are frozen dataclasses.

\begin{lstlisting}
@dataclass(frozen=True)
class CEdge:
    a: tuple        # loop-momentum routing coefficients (over the L loops)
    b: tuple        # external-q routing coefficients (over the n_ext externals)
    kind: str       # 'R' (retarded segment) or 'C' (correlation line)
    u: int          # endpoint vertex id
    v: int          # endpoint vertex id  (u == v  =>  self-loop = tadpole)
    external: bool  # True iff the edge touches an external leaf
    fpairs: tuple = ()   # coupled-field propagator matrix indices (Dyson 3c)
\end{lstlisting}

\begin{lstlisting}
@dataclass(frozen=True)
class CStackDiagram:
    internal_vertices: tuple   # interaction-vertex ids (their times integrated)
    external_legs: tuple       # leaf vertex ids, in correlation-function order
    edges: tuple               # tuple of CEdge
    n_loops: int               # == diagram loop_number
\end{lstlisting}

The mapping itself, reading \file{diagram\_descriptor.py:105-193}:

\begin{enumerate}
  \item Unpack \code{td.prediagram} into the graph \code{D}, the leaves, and the
        internal vertices; read \code{td.vertex\_assignments}.
  \item \textbf{Classify each non-leaf vertex.} A \code{VertexType} is an
        interaction vertex (its time is integrated). A 2-point \code{SourceType}
        --- the Gaussian noise insertion $\avg{\tilde\phi\tilde\phi}$ that
        represents a $C$ line --- is a \emph{noise} vertex, to be contracted away.
        A \code{SourceType} with $\ge 3$ response legs (non-Gaussian noise) is
        \emph{kept} as an internal vertex with $n$ outgoing $R$ edges --- the
        paper's all-retarded formulation --- and its time is integrated like any
        vertex. A source with fewer than 2 response legs raises
        \code{NotImplementedError}.
  \item Call \code{route\_momenta(td)} and read
        \code{ec = rr.edge\_coeffs()} into a per-edge $(a,b)$ map.
  \item Build an incidence map from each vertex to its incident edges.
  \item \textbf{Contract each 2-point noise source} into a single $C$ edge: its
        two incident $\Gret$ half-edges become one $C$ edge between the two
        \emph{other} endpoints (the source time integrated out gives
        $C(\Delta t)$). A noise source whose two edges land on the \emph{same}
        vertex becomes a $C$ \emph{self-loop} ($u=v$) --- the tadpole. Degree
        $\ne 2$ raises.
  \item \textbf{Every remaining $\Gret$ edge becomes an $R$ edge.} Edges touching
        a leaf are flagged \code{external=True}.
  \item Assemble and return the \code{CStackDiagram}, internal vertices sorted,
        external legs in correlation-function order, with
        \code{n\_loops = rr.n\_loops}.
\end{enumerate}

\begin{note}
\textbf{Why the $C$ edge may take either half's routing coefficients.} When the
noise source is contracted, the new $C$ edge takes the $(a,b)$ of one of its two
halves (\file{diagram\_descriptor.py:168}). This is correct because flipping
$(a,b)\to(-a,-b)$ leaves every Symanzik sum $\sum_e w_e a a$, $\sum_e w_e a b$,
$\sum_e w_e b b$ invariant --- the routing sign is irrelevant to the quadratic
forms. The descriptor \emph{relies} on this; do not try to ``fix'' the sign.
\end{note}

\begin{gotcha}
\textbf{There is no bubble / tadpole / sunset branch in the descriptor.} Whether a
loop momentum couples to the external $q$ (a bubble) or not (a tadpole) is a
property of the Symanzik $\Fsym$ \emph{downstream}, not of this descriptor
(\file{diagram\_descriptor.py:10-13}). \code{CStackDiagram.is\_tadpole\_like()}
exists (\file{diagram\_descriptor.py:93}) but is \emph{diagnostic only} --- the
production evaluator never branches on it. This topology-blindness is the whole
point: one integral, every diagram.
\end{gotcha}

\section{The kinematic integral: \texorpdfstring{\texttt{diagram\_kinematic}}{diagram\_kinematic}}\label{sec:kinematic}

\code{diagram\_kinematic} (\file{full\_integrator.py:469}) is the heart of the
subsystem. It computes the kinematic (unit-amplitude, no couplings) full-diagram
integral
$\int\prod\dd t_v\prod\dd\sigma_e\,\mathbf{1}(\theta)\,\ee^{-\mu\sum w}\,\text{MomFactor}$
by causal-chamber quadrature. Its signature is:

\begin{lstlisting}
def diagram_kinematic(descr, q_vec, external_times, mu, D, spatial_dim=1,
                      W=None, n_t=22, n_s=24, formfactor=None, gh_order=6,
                      xs=None, method='grid', mc_n=1000000, mc_seed=0):
\end{lstlisting}

Walking the body:

\begin{enumerate}
  \item \textbf{Extract the structure} (\file{full\_integrator.py:487-493}). Pull
        \code{edges}, \code{internal} vertices, the $a$/$b$ coefficient matrices
        (one row per edge), and $n_C$ (the number of $C$ edges).
  \item \textbf{Set the time window} (\file{full\_integrator.py:495-499}).
        $[\,\text{lo},\text{hi}\,]$ comes from the external times and a window
        $W=22/\mu$ --- about twenty-two relaxation times, well past where the
        integrand has decayed.
  \item \textbf{Build the retarded poset on the internal vertices}
        (\file{full\_integrator.py:503-515}). An internal$\to$internal $R$ edge
        becomes a poset edge \code{internal\_R}; an $R$ edge to or from a leaf
        becomes a fixed-time scalar bound \code{s\_up}/\code{s\_lo} (the leaf time
        clamps the vertex from one side).
  \item \textbf{Build the correlation Schwinger grid}
        (\file{full\_integrator.py:517-536}). This step is subtle enough to
        deserve its own paragraph below.
  \item \textbf{Dispatch on \code{method}} (\file{full\_integrator.py:544-555}).
        \code{'bessel'} routes to \code{\_diagram\_bessel\_xs}; otherwise
        enumerate \code{chambers = causal\_chambers(n\_V, internal\_R)}.
  \item \textbf{Per chamber, build the internal-time grid.} \code{'mc'}
        importance-samples the chamber via nested $\mathrm{Exp}(\mu)$ time gaps
        (\file{full\_integrator.py:564-578}); \code{'grid'} (the default) builds
        the nested causal-chamber product grid latest$\to$earliest, each level
        bounded above by the next-later time and its external scalar bound, via
        \code{\_gl\_on} (\file{full\_integrator.py:586-603}).
  \item \textbf{Compute the edge weights} \code{w\_batch} and the residual mass
        \code{mu\_resid} ($R$: $t_v-t_u$; $C$: $|t_u-t_v|+\sigma$) and the
        per-sample amplitude \code{amp} (\file{full\_integrator.py:606-623}).
  \item \textbf{Accumulate the result}, in one of four sub-cases
        (\file{full\_integrator.py:625-709}): a plain analytic-IFT heat kernel; a
        multivariate ($k\ge3$) derivative-vertex IFT; a single-external
        ($k=2$) derivative-vertex IFT; or the $q$-evaluation path. The last three
        are the subject of Chapter~\ref{ch:spatial-heatkernel}.
  \item \textbf{Return} a real vector for the analytic-IFT path, complex for a
        derivative vertex, float otherwise
        (\file{full\_integrator.py:714-716}).
\end{enumerate}

\subsection{The retarded-time quadrature: \texorpdfstring{\texttt{\_gl\_on}}{\_gl\_on}}

The internal vertex times are integrated by \code{\_gl\_on}
(\file{full\_integrator.py:349}), Gauss--Legendre nodes on $[\text{lo},\text{hi}]$
with a $\sqrt{\,}$-concentrating substitution toward the upper bound, where the
retarded integrand peaks:

\begin{lstlisting}
def _gl_on(lower, upper, n):
    xg, wg = np.polynomial.legendre.leggauss(n)
    v = 0.5 * (xg + 1.0)
    span = (upper - lower)[..., None]
    t = upper[..., None] - span * (v * v)   # nodes cluster near `upper`
    w = span * (wg * v)                     # Jacobian 2.span.v . 1/2
    return t, w
\end{lstlisting}

The $t = \text{upper} - \text{span}\cdot v^2$ map clusters nodes near the upper
bound (because $v^2$ is dense near $v=0$) and carries the Jacobian
$\text{span}\cdot(w_g\,v)$. This concentration matters: the retarded propagator
$\ee^{-m\,w}$ is largest when the segment length $w$ is smallest, i.e.\ when the
vertex time sits just below its upper bound.

\subsection{The Schwinger \texorpdfstring{$\sigma$}{sigma} grid is load-bearing}

The correlation Schwinger parameters $\sigma_e\in[0,\infty)$ are integrated by the
\emph{same} kind of substitution. Reading
\file{full\_integrator.py:517-525}:

\begin{lstlisting}
s_cap = 32.0 / mu
xv, wv = np.polynomial.legendre.leggauss(n_s)
vv = 0.5 * (xv + 1.0)
s_nodes = s_cap * vv * vv                         # sigma = s_cap . v^2
s_w = wv * s_cap * vv * np.exp(-mu * s_nodes)     # w.s_cap.v.e^{-mu.sigma}
\end{lstlisting}

That is Gauss--Legendre on a finite window $[0,s_{\text{cap}}]$ with
$\sigma = s_{\text{cap}}\,v^2$ and the $\ee^{-\mu\sigma}$ weight folded into the
node weights, \emph{not} Gauss--Laguerre on $[0,\infty)$.

\begin{gotcha}
\textbf{The $\sigma$-node concentration is not optional.} The substitution
$\sigma = s_{\text{cap}}\,v^2$ concentrates nodes near $\sigma = 0$
specifically to resolve the integrable self-loop singularity
$\Usym^{-d/2}\sim\sigma^{-d/2}$ (a $C$ self-loop has $\Usym = \sigma$, so as
$\sigma\to0$ the integrand has an integrable spike). Plain Gauss--Laguerre
\emph{under-resolves} that spike. The code comment at
\file{full\_integrator.py:519} says exactly this. Changing this quadrature will
quietly degrade tadpole accuracy without raising any error. (Note that the
function's own docstring at \file{full\_integrator.py:482} still calls this
``Gauss--Laguerre''; the body is the Legendre-plus-substitution above, and the
body is what runs.)
\end{gotcha}

\subsection{A worked shape: the one-loop bubble}

Make the abstract concrete with a one-loop bubble, $k=2$, $d=1$. It has two
internal vertices ($n_V=2$), two correlation lines ($n_C=2$), and one loop
($L=1$). \code{causal\_chambers(2, internal\_R)} returns the orderings of the two
vertices. Per chamber the grid has $P \approx 22^2\cdot 24^2$ samples; $a$ is
$(E,1)$ and $b$ is $(E,1)$; $\Lambda$ is $(P,1,1)$ so $\Usym = w_1+w_2$ and
$Q_{\text{eff}} = w_1w_2/(w_1+w_2)$; $\mathcal{B} = D\cdot Q_{\text{eff}}$ is a
scalar; and the heat kernel $(4\pi\mathcal{B})^{-1/2}\ee^{-x^2/4\mathcal{B}}$ is
summed over the grid into $\delta C(x)$ (the analytic IFT of
Chapter~\ref{ch:spatial-heatkernel}).

\section{The batched Symanzik core}\label{sec:batched}

The production integrator does \emph{not} call the single-sample
\file{spatial\_reduce.py} functions in its inner loop. Instead it builds the
Symanzik matrices for an entire grid of Schwinger samples at once, using NumPy's
Einstein-summation engine \code{np.einsum}.

\begin{note}
\textbf{What \code{np.einsum} is.} \code{np.einsum} evaluates a tensor
contraction written in index notation. The string \code{'pe,ej,ek->pjk'} reads
``take an array indexed by $(p,e)$ and two indexed by $(e,j)$ and $(e,k)$, sum over
the repeated index $e$, and keep $(p,j,k)$''. It runs in compiled C, and for the
heavy linear algebra (\code{det}, \code{solve}) NumPy calls BLAS/LAPACK, which
releases the Python GIL. Here $p$ indexes the Schwinger sample, $e$ the edge, $l/m$
the loop momenta, and $j/k$ the external momenta.
\end{note}

The batched factory is \code{\_momentum\_factor\_batch}
(\file{full\_integrator.py:56}). Its core (\file{full\_integrator.py:73-90}):

\begin{lstlisting}
Q   = np.einsum('pe,ej,ek->pjk', w_batch, b, b)   # (P, n_ext, n_ext)
if L == 0:
    quad = np.einsum('j,pjk,k->p', qv, Q, qv)
    return np.exp(-D * quad)                       # tree: exp(-D q^T Q q)
Lam = np.einsum('pe,el,em->plm', w_batch, a, a)    # (P, L, L)
N   = np.einsum('pe,el,ej->plj', w_batch, a, b)    # (P, L, n_ext)
U   = np.linalg.det(Lam)                           # (P,)  the first Symanzik
ok  = U > u_floor                                  # degeneracy mask
LamiN = np.linalg.solve(Lam[ok], N[ok])            # Lam^{-1} N, batched
Qeff  = Q[ok] - np.einsum('plj,plk->pjk', N[ok], LamiN)
pref  = (4.0*math.pi*D)**(-0.5*spatial_dim*L) * np.power(U[ok], -0.5*spatial_dim)
out[ok] = pref * np.exp(-D * np.einsum('j,pjk,k->p', qv, Qeff, qv))
\end{lstlisting}

Every line corresponds to one symbol in \eqref{eq:collapse}: \code{Q}, \code{Lam},
and \code{N} are the three Symanzik matrices for the whole batch; \code{U} is
$\det\Lambda$ for every sample at once; \code{LamiN} is $\Lambda^{-1}N$ via the
batched solve; \code{Qeff} is the Schur complement; and \code{pref} is the
$(4\pi D)^{-Ld/2}\Usym^{-d/2}$ prefactor.

\begin{gotcha}
\textbf{Degeneracy floors are everywhere, and they silently zero samples.} The
mask \code{ok = U > u\_floor} (with \code{u\_floor = 1e-300}) selects the samples
where the loop matrix is non-degenerate; degenerate samples (a self-loop where
$\Usym\to0$) get \emph{zero}, with no warning. The same pattern appears for the
heat-kernel width ($\mathcal{B} > 10^{-300}$, $\det\mathcal{B} > 10^{-300}$). A
diagram whose \emph{entire} grid is degenerate returns $0$ silently. This is
correct (those configurations have measure zero), but it means a bug that
degenerates every sample looks like a clean zero, not an error.
\end{gotcha}

The sibling \code{\_symanzik\_kernel\_batch} (\file{full\_integrator.py:94}) is the
analytic-IFT version: it returns the $q$-Gaussian
$\text{pref}\cdot\ee^{-q^{\mathsf T}\mathcal{B}q}$ \emph{without} contracting in $q$,
so the $x$-dependence stays analytic. That function feeds the heat-kernel IFT of
Chapter~\ref{ch:spatial-heatkernel}.

\section{The three backends}\label{sec:backends}

\code{diagram\_kinematic} takes a \code{method} argument selecting one of three
quadrature backends. They trade memory against generality.

\begin{description}
  \item[\code{'grid'} (default).] The deterministic causal-chamber product grid
        described in \S\ref{sec:kinematic}. Exact (to quadrature order), validated,
        but its grid size is $P \approx n_t^{\,n_V}\cdot n_s^{\,n_C}$ \emph{per
        chamber} --- the curse of dimensionality.
  \item[\code{'mc'}.] Monte-Carlo importance sampling
        (\file{full\_integrator.py:564-578}): internal times via nested
        $\mathrm{Exp}(\mu)$ gaps, correlation $\sigma$'s as $\mathrm{Exp}(\mu)$,
        with $P=$\code{mc\_n} random points instead of the product grid. Bounded
        memory, error $O(1/\sqrt{N})$.
  \item[\code{'bessel'}.] The radial-Bessel $\times$ angular-Monte-Carlo backend
        (\file{full\_integrator.py:362}), covered in detail in
        Chapter~\ref{ch:spatial-heatkernel}. It reparametrizes the causal-region
        point as $\lambda\cdot\hat s$ with $\hat s$ on a simplex, does the radial
        $\lambda$-integral \emph{exactly} as a modified Bessel function $K_\nu$,
        and samples only the smooth angular simplex. The radial reduction handles
        the $\det\Lambda\to0$ degenerate-loop direction analytically.
\end{description}

\begin{gotcha}
\textbf{The memory wall at $\ell\ge2$.} The grid is $P=n_t^{\,n_V}\cdot n_s^{\,n_C}$
per chamber. A KPZ two-loop diagram ($n_V=4$, $n_C=3$) is $\approx 1.8\times10^8$
points per chamber --- tens of gigabytes. The bridge refuses up front (raising a
\code{SpatialPropagatorError}) and points you to the escape hatches: lower
\code{max\_ell}, set \code{SPATIAL\_INTEGRATOR=mc} or \code{=bessel}, coarsen
\code{SPATIAL\_GRID\_NT}/\code{NS}, or raise \code{SPATIAL\_MEM\_BUDGET\_GB}.
\end{gotcha}

\begin{gotcha}
\textbf{Monte-Carlo is biased for derivative vertices.} The \code{'mc'} backend is
validated to better than $0.1\%$ for \emph{plain} $\phi^n$ vertices but is
\emph{biased} for derivative (form-factor) vertices: the $\det\Lambda\to0$
singularity gives the Monte-Carlo estimator infinite variance. For derivative
vertices at $\ell\ge2$, use \code{SPATIAL\_INTEGRATOR=bessel} --- the radial
reduction cures the variance. This is stated repeatedly in the source
(\file{full\_integrator.py:560-563}).
\end{gotcha}

\section{Assembling a correlator}

The kinematic integral is unit-amplitude. The amplitude, the retarded$+$advanced
completion, and the diagram sum are layered on by a short chain of functions, all
near the bottom of \file{full\_integrator.py}.

\code{diagram\_value} (\file{full\_integrator.py:1142}) applies the amplitude:
\begin{lstlisting}
def diagram_value(descr, prefactor_val, q_vec, external_times, mu, D,
                  spatial_dim=1, **kw):
    n_C = sum(1 for e in descr.edges if e.kind == 'C')
    kin = diagram_kinematic(descr, q_vec, external_times, mu, D,
                            spatial_dim=spatial_dim, **kw)
    return (2.0 ** (-n_C)) * float(prefactor_val) * kin
\end{lstlisting}
That is exactly $2^{-n_C}\cdot\Sym(\Gamma)\cdot\text{kinematic}$, with the
$2^{-n_C}$ of \S\ref{sec:norm}.

\code{diagram\_correlator} (\file{full\_integrator.py:1166}) adds the
retarded$+$advanced completion. A \emph{retarded-type} insertion --- one whose two
external legs have \emph{different} propagator kinds, one $C$ and one $R$
(\code{\_is\_retarded\_type}, \file{full\_integrator.py:1156}) --- dresses both the
retarded and advanced sides of the line, so it appears as the pair
$\Sigma_R(\tau)+\Sigma_A(\tau) = \Gamma(\tau)+\Gamma(-\tau)$. A $\{R,R\}$ (Keldysh)
insertion is its own conjugate and contributes a single $\Gamma(\tau)$. At
$\tau=0$ the retarded-type case becomes $2\Gamma(0)$.

\code{correlator\_2pt} (\file{full\_integrator.py:1182}) sums over all enumerated
diagrams --- tree plus every loop, each via the \emph{same} full integral ---
skipping diagrams whose prefactor is numerically dead
($|\text{pre}| < 10^{-14}$, ``dead at the saddle''):
\begin{lstlisting}
def correlator_2pt(descrs_prefactors, q, tau, mu, D, spatial_dim=1, **kw):
    total = 0.0
    for descr, pre in descrs_prefactors:
        if abs(float(pre)) < 1e-14:
            continue
        total += diagram_correlator(descr, pre, q, tau, mu, D,
                                    spatial_dim=spatial_dim, **kw)
    return total
\end{lstlisting}

The three functions \code{diagram\_value\_x}, \code{diagram\_correlator\_x}, and
\code{correlator\_2pt\_x} (\file{full\_integrator.py:1203},
\file{full\_integrator.py:1213}, \file{full\_integrator.py:1228}) are the
analytic-IFT analogues: they return real-space $\delta C(x,\tau)$ as a vector over
the output positions \code{xs}, directly, with no $q$-grid. \textbf{These are the
production entry points} the bridge calls for real-space output, and they are the
bridge to Chapter~\ref{ch:spatial-heatkernel}.

\begin{gotcha}
\textbf{Complex per-diagram values are normal.} A derivative vertex with an odd
number of derivatives gives a complex per-diagram contribution, because
$\partial_x\to\ii k$ in momentum space. The \emph{physical} correlator is real:
the imaginary parts cancel in the diagram sum and are dropped at the real-space
output (\file{full\_integrator.py:710-716}). Do not panic at a complex
intermediate value --- it is a bookkeeping artifact of a single diagram, not a
sign of error.
\end{gotcha}

\section{The oracle evaluator}

The sixth file, \file{generic\_evaluator.py}, is \emph{not} on the production
path. Its banner says so in the first lines (\file{generic\_evaluator.py:2-4}):
``ORACLE-ONLY --- not on the production path. Superseded by
\code{full\_integrator.py} $\dots$ reached only by its own test(s). Kept as an
independent numerical cross-check --- \code{compute\_cumulants} does NOT use this
module.'' It evaluates a diagram by an older \emph{Dyson-convolution} route
(self-energy $\sigma(q,u)$ followed by a single-mode Dyson dressing).

\begin{gotcha}
\textbf{The oracle still branches on topology --- which is exactly what the
production design avoids.} \file{generic\_evaluator.py} branches on
\code{is\_tadpole\_like()} to split the bubble and tadpole cases (in its
\code{diagram\_delta\_C} dispatch). That is the \emph{only} place in the spatial
subsystem where a topology branch survives, and it is a relic of the older design.
It is kept solely so the topology-blind production integrator can be cross-checked
against an independent implementation. Never call it from production code.
\end{gotcha}

\section{What feeds and consumes this subsystem}

To close the loop on the data flow: the orchestrator that prepares the inputs and
consumes the outputs is \file{msrjd/integration/spatial/pipeline\_bridge.py}.

\textbf{In.} \code{build\_pipeline\_records} yields $(td,\text{pre})$ pairs per
loop order --- a \code{TypedDiagram} and a symbolic prefactor. The bridge
substitutes numeric couplings to get a float
$\text{pv} = \text{float}(\text{pre}\big|_{\text{saddle}})$, builds the descriptor
\code{dd = diagram\_to\_cstack(td)}, optionally a form factor \code{ff}, and tags
the loop order. Diagrams with $|\text{pv}| < 10^{-14}$ are dropped.

\textbf{Through.} For real-space output the bridge calls, per live diagram and per
$\tau$ (\file{pipeline\_bridge.py:1697}):
\begin{lstlisting}
dCx_by_ell[el][it, :] += diagram_correlator_x(
    dd, pv, xg, float(tau), mu0, D0, spatial_dim=d,
    n_t=nt, n_s=ns, formfactor=ff, ...)
\end{lstlisting}
which flows \code{dd, pv} $\to$ \code{diagram\_value\_x} $\to$
\code{diagram\_kinematic(..., xs=xg)} $\to$ the chamber loop $\to$
\code{\_symanzik\_kernel\_batch} $\to$ \code{\_heat\_kernel\_x\_general} $\to$ the
accumulated $\delta C(x)$.

\textbf{Out.} The bridge accumulates an array of shape $(n_\tau, n_x)$ per loop
order, which becomes \code{compute\_cumulants}'s $\delta C(x,\tau)$.

\begin{note}
\textbf{macOS fork crash --- a project-wide rule that touches this subsystem only
by mandate.} Fork-based multiprocessing crashes the user's Jupyter kernel and
operating system on this machine; the spatial path is therefore serial or
thread-only by mandate. \emph{None} of the six files in this subsystem fork --- but
if you extend the integrator, any parallelism you add must be thread-based (NumPy
releases the GIL for the heavy linear algebra) or $q$-grid batching, never
process- or fork-based.
\end{note}

\section{Recap}

You have now seen the spine of the spatial integrator:
\begin{enumerate}
  \item \textbf{Why momentum-first} --- the close-pair pathology lives in the
        momentum integral, so doing that integral analytically removes it
        (\S\ref{sec:heatkernel}).
  \item \textbf{Momentum routing} --- conservation at each spatial vertex, solved
        symbolically with SymPy, gives the routing coefficients $(a_e,b_e)$
        (\S\ref{sec:routing}).
  \item \textbf{Symanzik polynomials} --- the loop Gaussian collapses to
        $(4\pi D)^{-Ld/2}\Usym^{-d/2}\ee^{-Dq^{\mathsf T}Q_{\text{eff}}q}$, with
        $\Usym=\det\Lambda$ and $Q_{\text{eff}}$ the Schur complement, exact in $L$
        and $d$ (\S\ref{sec:symanzik}).
  \item \textbf{The descriptor} --- \code{diagram\_to\_cstack} flattens every
        typed diagram into one canonical form with no topology branch
        (\S\ref{sec:descriptor}).
  \item \textbf{Causal chambers} --- the retarded poset's linear extensions tile
        the time domain into smooth chambers (\S\ref{sec:chambers}).
  \item \textbf{The kinematic integral} --- \code{diagram\_kinematic} does the
        residual smooth integral by batched Gauss--Legendre product grids over the
        chambers (\S\ref{sec:kinematic}, \S\ref{sec:batched}), with three backends
        trading memory for generality (\S\ref{sec:backends}).
\end{enumerate}

What this chapter \emph{deferred} is the step from the per-sample $q$-Gaussian to
real-space $\delta C(x,\tau)$ --- the analytic inverse Fourier transform --- and
the handling of derivative (form-factor) vertices. Both are
Chapter~\ref{ch:spatial-heatkernel}. The coupled multi-field (unequal-$D$)
extension, which replaces the single mass $\mu$ with a table of per-segment masses,
is Chapter~\ref{ch:spatial-coupled}.
